\documentclass[instructions.tex]{subfiles}
\begin{document}
 
\chapter{De la comédie et des spectacles.}
 
\primo Si la comédie se bornait à représenter avec décence des exemples édifians, ou les actions mémorables des grands hommes, elle ne serait point condamnable; mais ce n’est point là ce qu’on y voit. Tout ce qui est capable de réveiller les passions, d’exciter la concupiscence de la chair et des yeux, et l’orgueil de la vie, s’y réunit. Car, sans parler du concours et des rendez-vous de la jeunesse de tout sexe, à qui la comédie est une occasion de désordre, jugeons de la comédie par les circonstances, et par les sujets qui y sont représentés.
 
1. Les circonstances et l’appareil de la comédie, les décorations agréables et enchantées, la vue des actrices, leurs parures, leur enjouement, leurs voix insinuantes, les airs tendres et passionnés des acteurs, les tours délicats sur la pudeur et l’amour profané; les traits satiriques lâchés, en passant, sur la vertu; tout cela ne fait-il aucune impression sur les cœurs? Si l’on a peine à résister à ces impressions, quand seul, y résistera-t-on dans la dissipation du spectacle?
 
2. Quant aux sujets qui sont le fond et la base de la comédie, sans compter les bouffonneries, les extravagances, les sauts et les gestes dissolus, ces femmes et ces acteurs qui exposent leur vie, en se balançant indécemment sur des cordes, que voit-on dans le reste, qu’une peinture des passions plus propre à les exciter, qu’à les éteindre? tantôt une intrigue de galanterie, une maîtresse affligée, un rival soupirant, une femme jalouse, un mari dupé; tantôt des satircs piquantes contre les différens états; d’autres fois des aventures tragiques, des trahisons, des fourberies, des combats, des vengeances méditées, des projets ambitieux exécutés avec succès, une conspiration, des cruautés exercées avec fureur; quelquefois même la religion, les personnes sacrées et les puissances tournées en ridicule, etc.
 
En vérité, un chrétien peut-il se croire innocent dans le plaisir qu’il prend à voir, à entendre ce qui excite en lui tant de passions différentes? Et quand il serait sans passion, lui est-il permis de voir avec complaisance les représentations des choses qu’il doit détester? Dieu, qui par la sainteté de sa loi nous ordonne de veiller en tout temps sur nos sens, sur notre esprit et sur notre cœur, pour en écarter les représentations et les pensées dangereuses, qui fera rendre compte d’une parole inutile, et des moindres dépenses superflues, peut-il approuver des spectacles qui remplissent l’esprit et l’imagination de tant d’objets vains, ridicules et séducteurs; peut-il approuver qu’on y emploie un argent dont on devrait soulager tant de pauvres familles qui gémissent dans l’indigence?
 
\secundo Le monde cependant prétend avoir de grandes raisons pour les autoriser. La comédie, dit-on, est utile, elle déclame contre le vice, attirant que les prédicateurs. Quelle indignité de mettre le théâtre en parallèle avec l’Evangile, et de comparer la parole d’un comédien avec la parole de Dieu! La comédie, il est vrai, rend le vice ridicule, mais elle ne le rend pas odieux; elle en fait rire, mais elle ne le fait pas pleurer; elle inspire la ruse, la défiance, le mépris d’autrui, la satire, non la charité; elle a fait commettre des millions de péchés, et jamais elle n’en a fait détester un seul.
 
Vous prêchez contre la comédie, me dit un jour un homme qui avait été parmi les acteurs sur un théâtre; vous avez raison; elle fait commettre cent fois plus d’ennemies que vous ne pouvez imaginer. Les fruits qui croissent sur les bords du lac de Sodome paraissent d’une beauté charmante; mais aussitôt qu’on les touche, ils tombent en poussière, et répandent une infection insupportable. Tels sont les fruits de la comédie; en s’évanouissant, ils répandent dans l’âme un air contagieux.
 
Mais, dira-t-on, je n’y vais que par divertissement; je n’y ai jamais eu ni mauvaises pensées, ni tentations.
 
Vous vous trompez. Étourdi par l’enchantement du spectacle, vous n’avez pas connu ce qui se passait en vous. Dans le saint lieu même, souvent vous avez eu des tentations: comment n’en auriez-vous pas à la comédie? Vous avez pensé, dans ces spectacles, aux objets que vous y voyiez et à ce qu’on y disait; vous en sortiez avec un esprit rempli d’idées profanes, qui vous ôtaient le goût des choses de Dieu et de vos devoirs; qui vous dissipaient jusque dans le temps de prier ce qui agitait en vous tout-à-tour différentes passions; ne sont-ce pas autant de tentations? S’il vous faut quelque divertissement, faites comme d’autres, qui, sans aller aux bals et à la comédie, savent se divertir innocemment.
 
Mais, ajoute-t-on, saint François de Sales ne condamne pas les danses et les spectacles. Cela est faux. Loin d’approuver ces sortes de divertissemens, il a écrit tout ce qui est capable d’en faire connaître le ridicule, le danger et le venin. Ce grand saint, à la vérité, en faveur de ceux qui, dans certaines conjonctures qui sont rares, se voient comme forcés de s’y trouver, prescrit des précautions pour y conserver l’innocence; mais il ne plaît pas aux gens du monde de les prendre, ces précautions; ils ont donc mauvaise grâce d’autoriser leurs danses et les spectacles par le témoignage du saint évêque.
 
N’alléguez point qu’étant lié avec le monde, vous ne pouvez vous dispenser de faire comme les autres, ni vous passer de ces sortes de divertissemens. Saint Augustin vous répondra que bien d’autres plus distingués et meilleurs que vous, s’en dispensent et s’en passent: pourquoi ne pourriez-vous pas faire de même \footnote{Numquid delicatior es illo senatore? tu non potes ? ille petuit. S.Aug.} ?
 
Il faut donc, ajoutez-vous, vivre comme des solitaires et des misanthropes. D’ailleurs ne vaut-il pas mieux aller à la comédie et au bal que de faire plus de mal?
 
Un pareil discours dans la bouche d’un chrétien est un raisonnement insensé, qui ne mérite pas qu’on y réponde \footnote{Ne respondeas stulto juxtà stultitaem suam.} .
 
\section{Résolutions}
 
\primo Je renonce pour toujours aux spectacles.

\secundo Et je ne permettrai pas à ceux qui seront à ma charge de les fréquenter.
 
\prayer{Mon Dieu, faites-moi la grâce de remplir exactement les promesses que je vous ai faites dans mon baptême, où j’ai abjuré satan, ses pompes et ses œuvres.}
 
\end{document}
